\section{System moments, initial conditions and the practical panel choice}
\label{s5:dynamic-panel}
The difference equation removes an individual effect; it does not create a
strong instrument. This distinction matters for persistent loan balances or
capital ratios. Consider first the scalar model
\[
 y_{it}=\alpha_i+\rho y_{i,t-1}+v_{it},\qquad |\rho|<1,
 \quad \E[v_{it}\mid\alpha_i,y_{i1},v_{i2},\ldots,v_{i,t-1}]=0.
\]
Independent entities and fixed $T$ define the sampling experiment. For
$s\le t-2$, both $\E[y_{is}v_{it}]$ and $\E[y_{is}v_{i,t-1}]$ vanish, giving
$\E[y_{is}\Delta v_{it}]=0$. These are the differenced moments already
derived above \citep{arellano1991}. Their validity does not require the
initial observation to have the stationary distribution.

Level equations permit a different restriction. Write
$y_{i1}=\alpha_i/(1-\rho)+e_{i1}$. Then
\[
 \Delta y_{i2}=(\rho-1)e_{i1}+v_{i2},\qquad
 \E[\Delta y_{i2}\alpha_i]=0
 \quad\text{if }\E[e_{i1}\alpha_i]=0.
\]
The recursion
$\Delta y_{it}=\rho^{t-2}\Delta y_{i2}+
\sum_{j=0}^{t-3}\rho^j\Delta v_{i,t-j}$ propagates this zero covariance.
Consequently $\E[\Delta y_{i,t-1}(\alpha_i+v_{it})]=0$.
Combining these level moments with differenced moments yields system GMM.
Mean stationarity is sufficient for this construction; equality of initial
and limiting variances is not required. Conversely a portfolio assembled at
loan origination or after a balance-sheet repair may systematically start
above or below its entity-specific long-run level. The additional moment can
then fail even when the differenced moments remain valid. This reconstructs
the initial-condition argument in \citet[Sections 3.1--3.3, pp.107--109]{blundellbond2023}.
The archived text is their 2023 account, not the original 1998 article.

Adding $\beta x_{it}$ requires a separate classification for $x$.
Contemporaneous endogeneity excludes $x_{it}$ as an instrument; lagged levels
must precede both innovations in $\Delta v_{it}$. Predetermination and strict
exogeneity allow different lag sets. A level moment using $\Delta x_{i,t-1}$
also needs $\E[\Delta x_{i,t-1}\alpha_i]=0$ and orthogonality to $v_{it}$.
No initial-condition theorem for $y$ automatically supplies this condition
for a bank's chosen action. With common macro shocks, cross-entity
independence itself fails and the fixed-$T$, large-$N$ argument must change.

\subsection{Weighting, covariance and what a correction can repair}
Let $b_i-G_i\theta$ collect the entity's moments after any prespecified lag
restriction and compression, where $\theta=(\rho,\beta)'$. For a positive
weight $W$, the quadratic criterion has minimizer
\[
 \hat\theta_W=(\bar G'W\bar G)^{-1}\bar G'W\bar b,
 \quad
 \sqrt N(\hat\theta_W-\theta_0)=
 (G'WG)^{-1}G'W\,N^{-1/2}\sum_i(b_i-G_i\theta_0)+o_p(1).
\]
The expansion needs $\bar G\to G$, a moment CLT, a nonsingular $G'WG$,
and consistent $W$. Thus the covariance estimator is
\[
 \widehat{\Var}(\hat\theta_W)=N^{-1}
 (\bar G'W\bar G)^{-1}\bar G'W\hat\Omega W\bar G
 (\bar G'W\bar G)^{-1},
\]
where $\hat\Omega$ is the covariance of the entire entity moment vector,
not of individual entity-date rows. Two-step GMM takes
$W=\hat\Omega(\hat\theta_1)^{-1}$ using an initial consistent estimator.
The first-order limit omits uncertainty in that estimated weight because
the population moment is zero. The finite-sample derivative is nevertheless
nonzero:
\[
 d\hat\theta_2=(\bar G'W\bar G)^{-1}\bar G'(dW)
 (\bar b-\bar G\hat\theta_2),\qquad
 dW=-W(d\hat\Omega)W.
\]
This displays the term whose variation a finite-sample correction must
address. The 2023 primary account explains the Windmeijer correction and its
motivation \citep[Section 2.3, p.105]{blundellbond2023}. The transparent local
reference implements the first-order sandwich, \emph{not} that correction;
the correction's original paper has not been substituted by a metadata record.
A production two-step adapter needs a source-qualified implementation of the
correction or a separately justified complete-refit alternative. Neither can
make invalid moments valid or rescue a rank-deficient Jacobian.

Compression also changes a statistical problem. An unrestricted triangular
lag matrix has $(T-2)(T-1)/2$ instruments in the scalar differenced model.
The reference collapses moments over dates and keeps lags two and three,
yielding four differenced moments for $(y,x)$, or six after adding level
moments. Scaling moments by observed root mean squares changes the one-step
weight; it leaves the population zero restrictions unchanged. The code
therefore reports this estimator as scaled identity-weight GMM, not the
canonical homoskedastic optimal one-step estimator. Small objective values
with many instruments can reflect overfitting. An overidentification test
is neither proof of exclusion nor evidence of adequate instrument strength.

\subsection{A paired comparison and a banking decision rule}
The synthetic short-panel experiment uses $N=800$, eight transitions after
an initial observation, and
$y_{it}=\alpha_i+\rho y_{i,t-1}+.6x_{it}+v_{it}$, with
$x_{it}=.5x_{i,t-1}+.4v_{it}+\eta_{it}$. A 400-period burn-in generates
the regular initial law. The initial-condition challenger then adds a
multiple of $\alpha_i$ to the initial balance; the serial-error challenger
changes the innovation law. The four methods use identical datasets.
Coverage below is joint for both coefficients, using two-coordinate
Bonferroni 95\% intervals across 400 repetitions.
\begin{center}\small
\begin{tabular}{lrrrr}\toprule
Regime & Difference 1 & Difference 2 & System 1 & System 2\\\midrule
Stationary & .955 & .950 & .958 & .948\\
Persistent & .985 & .983 & .968 & .933\\
Initial departure & .950 & .960 & .033 & .000\\
Serial-error violation & .000 & .000 & .000 & .000\\\bottomrule
\end{tabular}\end{center}
One and two denote weighting steps, not different causal targets. Near
95\% coverage, the Monte Carlo standard error is about .011. The very high
coverage of a weak differenced design is not a precision advantage; weak
moments can produce broad intervals and unreliable normal approximations.
The initial-departure failure is the decisive exclusion from an automatic
system-GMM default. For established comparable entities, system moments can
be a challenger after a substantive initial-law argument. For new portfolios,
the primary differenced route has the weaker initial-law requirement, but
must still report Jacobian strength, instrument count and residual serial
correlation. If those fail, the result is an unsupported inferential claim,
not a prompt to search until a small standard error appears.

Misspecification-robust moment inference is another distinct alternative.
\citet[Sections 2--3, Assumptions 1--2]{kleibergen2021} targets the minimizer
of a continuously updated population GMM criterion when no parameter makes
all moments zero. Its joint moment/Jacobian CLT and covariance conditions
support an identification- and misspecification-robust score construction.
That pseudo-true minimizer need not equal a causal response or a meaningful
risk premium. This use of ``double robust'' concerns weak identification and
moment misspecification; it is different from outcome/propensity robustness
in Chapter~\ref{r3:learning}. The inspected copy is the 2021 preprint of work
identified in 2025 journal metadata. It is a substantive inference challenger,
not a license to keep a causal label after the identifying restrictions fail.
